% part: intuitionistic-logic
% chapter: propositions-as-types
% section: type-preservation

\documentclass[../../../include/open-logic-section]{subfiles}

\begin{document}

\olfileid{int}{pty}{tpr}

\olsection{Type Preservation}

The reduction and permutation conversions considered can also be
defined directly for the systems \Log{tN2i} and \Log{tN3i}. The system
\Log{tN3i} defines the type of a $\lambd$-term~$N$ relative to the
context~$\Gamma$ as the !!{formula}~$!A$ such that $\Gamma \Sequent N
: !A$ has !!a{proof} in \Log{tN3i}. The fact that the
$\beta$-reduction rules for $\lambd$-terms exactly mirror the
reduction and permutation conversions of !!{proof}s has an important
consequence: the type of a $\lambd$-term relative to a
context~$\Gamma$ remains the same after $\beta$-conversion. This
property is called \emph{type preservation}.

\begin{prop}
If $N$ has type~$!A$ relative
to $\Gamma$ and $N \redone N'$, then $N'$ has type~$!A$ relative to~$\Gamma$.
\end{prop}

\begin{proof}
  By induction on~$N$ and the height of !!{proof}s~$\delta$ of $\Gamma
  \Sequent N : !A$ we can easily establish the following: If $M$ is a
  sub-term of~$N$, then $\delta$ contains a sub-!!{proof}~$\delta_1$
  ending in $\Gamma' \Sequent M : !B$. If $M$ is a redex, then a
  conversion applies to~$\delta_1$. By applying the reduction to the
  $\delta_1$ we obtain !!a{proof}~$\delta_1'$ of $\Gamma' \Sequent M'
  : !B$. By inspection of the conversions of \Log{tN3i}, $M'$ is the
  reductum of~$M$. Replacing $\delta_1$ in~$\delta$ by $\delta_1'$
  results in !!a{proof}~$\delta'$ of $\Gamma \Sequent N' : !A$, where
  $N'$ results from~$N$ by replacing the subterm~$M$ by~$M'$.

  For instance, consider the reduction conversion for~\Intro{\land}
  followed by \Elim{\land}:
  \[
  \bottomAlignProof
  \AxiomC{}
  \Deduce$\Gamma \fCenter N_1 : !B_1$
  \AxiomC{}
  \Deduce$\Gamma \fCenter N_2 : !B_2$
  \RightLabel{\Intro{\land}}
  \BinaryInf$\Gamma \fCenter \pair{N_1}{N_2} : !B_1 \land !B_2$
  \RightLabel{\Elim{\land}[i]}
  \UnaryInf$\Gamma \fCenter \proj{i}{\pair{N_1}{N_2}} : !B_i$
  \DisplayProof
  \quad\redone\quad
  \bottomAlignProof
  \AxiomC{}
  \Deduce$\Gamma \fCenter N_i : !B_i$
  \DisplayProof
  \]
  We see that the proof term of the !!{proof} on the right, i.e.,
  $\delta_1'$ is $N_i$. It is exactly the result of applying
  $\beta$-conversion to the proof term $\proj{i}{\pair{N_1}{N_2}}$ of
  the !!{proof} on the left~($\delta_1$).

  Or consider a \Intro{\lif} followed by an \Elim{\lif} inference, and
  the reduction conversion applied to it:
  \begin{multline*}
  \bottomAlignProof
  \AxiomC{$\delta_2$}
  \Deduce$x: !B_1, \Gamma \fCenter N : !B_2$
  \RightLabel{\Intro{\lif}}
  \UnaryInf$\Gamma \fCenter \lambd[\typeof{x}{!B_1}][N] : !B_1 \lif !B_2$
  \AxiomC{}
  \RightLabel{$\delta_3$}
  \Deduce$\Gamma \fCenter M : !B_1$
  \RightLabel{\Elim{\lif}}
  \BinaryInf$\Gamma \fCenter (\lambd[\typeof{x}{!B_1}][N] M) : !B_2$
  \DisplayProof
  \quad\redone\\
  \bottomAlignProof
  \AxiomC{}
  \RightLabel{$\delta_3$}
  \Deduce$\Gamma \fCenter M : !B_1$
  \RightLabel{$\Subst{\delta_2}{\delta_3}{x:!B_1}$}
  \Deduce$\Gamma \fCenter \Subst{N}{M}{x} : !B_2$
  \DisplayProof
  \end{multline*}
  Again the proof term of the !!{proof} on the right is exactly the
  result of $\beta$-reducing the proof term of the !!{proof} on the
  left. Of course one would have to carefuly verify (by induction on
  the height of~$\delta_2$) that if we replace all axioms of the form
  $\Gamma' \Sequent x:!B_1$ in~$\delta_2$ by the !!{proof}~$\delta_3$
  of $\Gamma \Sequent M:!B_1$ actually results in !!a{proof}
  $\Subst{\delta_2}{\delta_3}{x:!B_1}$ of $\Gamma \Sequent
  \Subst{N}{M}{x} : !B_2$.
\end{proof}

The tight correspondence between !!{proof}s and $\lambd$-terms has
another consequence. If $N$ is a $\lambd$-term of type $!A$ relative
to a context~$\Gamma$ and contains a redex~$M$ (i.e., it is not a
normal form), then the corresponding \Log{tN3i} proof~$\delta$ of
$\Gamma \Sequent N : !A$ contains a sub-!!{proof} (ending in $\Gamma'
\Sequent M: !B$) to which a conversion can be applied. If $N \redone
N'$ by replacing $M$ by its reductum, the corresponding result of
applying the reduction to~$\delta$ is !!a{proof} of $\Gamma \Sequent
N' : !A$. Hence every $\lambd$-term not in normal form $\beta$-reduces
to another. This property is called \emph{progress}.

Progress together with type preservation ensure that $\beta$-reduction
of $\lambd$-terms never ``gets stuck'': if a term is well-typed and
isn't in normal form, it reduces to another well-typed term (and that
term is of the same type).

\end{document}
